% !PRINTED 13
\weberpage{13}
%<<CONT
of two ratios with the same numerator and unequal denominators, that
one is smaller whose denominator is greater. For if $a,a',b$ are
arbitrary elements of a measurable set and $a'>a$, choose first a whole
number $n$ such that $na>b$ and $n(a'-a)>b$. Then take $m$ to be the
least whole number satisfying $mb>na$; then $na<mb$, but $mb<na'$.
% [source-original sic] The print gives mb\ge na'\ge na+n(a'-a); the second relation is algebraically equality and is preserved verbatim.
For if $mb\ge na'\ge na+n(a'-a)$, then $mb>na+b$; hence, contrary to
the supposition, already $(m-1)b>na$. Thus
\[
 {a\over b}<{m\over n}<{a'\over b},
\]
and hence $a:b<a':b$. In exactly the same way one proves that if
$b'>b$, then $a:b>a:b'$. This theorem is also expressed by saying that
a ratio increases with the numerator and decreases with increasing
denominator.

From this one easily obtains the following theorem. If $a,b,c,d$ are
elements of the same measurable set and $a:b=c:d$, then also
$a:c=b:d$. For suppose that $a:c<b:d$; then there would have to be two
whole numbers $m,n$ such that
\[
 na<mc,
 \qquad nb>md.
\]
Then, however, by the theorem just proved, $na:nb<mc:md$, and hence
$a:b<c:d$, contrary to the supposition.

\gesperrt{To this is now attached the following principal theorem. If
of four quantities} $a,b,c,d$ \gesperrt{any three are given arbitrarily
from a continuous measurable set, then the fourth can be determined in
the same set so that} $a:b=c:d$.

The theorem is an immediate consequence of the assumed continuity. For
if, for example, an element $x$ of $\mathfrak M$ is put into $A$ or
into $B$ according as $x:b$ is smaller or greater than $c:d$, then a
cut is obtained, which is produced by an element $a$ satisfying
$a:b=c:d$. But this theorem also holds in certain non-continuous sets,
for example for the rational fractions.

It follows that as representatives of two numbers one can always
choose two ratios whose elements belong to the same
%CONT>>
